\chapter{Kalman Score Recursions}
\label{ch:bf-kalman-score}

The score recursion differentiates the prediction-error decomposition from
Chapter~\ref{ch:bf-prediction-error-decomposition}.  In BayesFilter this is the
first derivative contract for exact linear Gaussian HMC targets.

For parameter component $\theta_i$, write
\[
  \dot z^{(i)} = \frac{\partial z}{\partial \theta_i}.
\]
All objects in the filtering recursion may depend on $\theta$:
\[
  c_t,\quad F_t,\quad Q_t,\quad a_t,\quad H_t,\quad R_t.
\]
Time subscripts are kept in this chapter because missing-data masks and
mixed-frequency observation systems can make the effective observation equation
time varying even when the underlying structural model is time homogeneous.

\section{Prediction Derivatives}
\label{sec:bf-score-prediction-derivatives}

The differentiated prediction step is
\begin{align}
\label{eq:bf-score-predicted-mean-first}
  \dot m_{t|t-1}^{(i)}
  &=
  \dot c_t^{(i)}
  +\dot F_t^{(i)}m_{t-1|t-1}
  +F_t\dot m_{t-1|t-1}^{(i)},\\
\label{eq:bf-score-predicted-cov-first}
  \dot P_{t|t-1}^{(i)}
  &=
  \dot F_t^{(i)}P_{t-1|t-1}F_t^\top
  +F_t\dot P_{t-1|t-1}^{(i)}F_t^\top
  +F_tP_{t-1|t-1}\dot F_t^{(i)\top}
  +\dot Q_t^{(i)}.
\end{align}
These equations also define the initial-condition contract: if the
implementation fixes $(m_{0|0},P_{0|0})$ as data, the initial derivatives are
zero; if it uses a stationary initial covariance, the derivatives must come from
the corresponding derivative Lyapunov equations and be reported as part of the
structural provider.

\section{Innovation Derivatives}
\label{sec:bf-score-innovation-derivatives}

The source note derives the innovation derivatives
\begin{align}
\label{eq:bf-score-innovation-first}
  \dot v_t^{(i)}
  =
  -\dot a_t^{(i)}
  -\dot H_t^{(i)}m_{t|t-1}
  -H_t\dot m_{t|t-1}^{(i)},\\
\label{eq:bf-score-innovation-cov-first}
  \dot S_t^{(i)}
  =
  \dot H_t^{(i)}P_{t|t-1}H_t^\top
  +H_t\dot P_{t|t-1}^{(i)}H_t^\top
  +H_tP_{t|t-1}\dot H_t^{(i)\top}
  +\dot R_t^{(i)}.
\end{align}
Using $\partial_i S^{-1}=-S^{-1}\dot S^{(i)}S^{-1}$, the per-time score
contribution is
\begin{equation}
\label{eq:bf-kalman-score-contribution}
  \frac{\partial \ell_t}{\partial \theta_i}
  =
  -\frac{1}{2}
  \left[
    \tr(S_t^{-1}\dot S_t^{(i)})
    +2\dot v_t^{(i)\top}S_t^{-1}v_t
    -v_t^\top S_t^{-1}\dot S_t^{(i)}S_t^{-1}v_t
  \right].
\end{equation}
For production code, BayesFilter prefers the solve form.  Let $S_tw_t=v_t$.
Then the same score can be written as
\begin{equation}
\label{eq:bf-solve-score}
  \frac{\partial \ell_t}{\partial \theta_i}
  =
  -\frac{1}{2}
  \left[
    \tr(S_t^{-1}\dot S_t^{(i)})
    +2\dot v_t^{(i)\top}w_t
    -w_t^\top\dot S_t^{(i)}w_t
  \right].
\end{equation}
The trace term may be computed by explicit materialization in small tests, but
production code should prefer factor solves or trace estimators appropriate to
the dimension and backend.

\section{Gain and Update Derivatives}
\label{sec:bf-score-update-derivatives}

The score contribution at time $t$ can be computed as soon as
\eqref{eq:bf-score-innovation-first}--\eqref{eq:bf-score-innovation-cov-first}
are available.  The recursion must still propagate derivatives to the next time
step.  For the covariance-form Kalman update,
\[
  K_t=P_{t|t-1}H_t^\top S_t^{-1},
\]
the gain derivative is
\begin{align}
\label{eq:bf-score-gain-first}
  \dot K_t^{(i)}
  &=
  \dot P_{t|t-1}^{(i)}H_t^\top S_t^{-1}
  +P_{t|t-1}\dot H_t^{(i)\top}S_t^{-1}
  -P_{t|t-1}H_t^\top S_t^{-1}\dot S_t^{(i)}S_t^{-1}.
\end{align}
Equivalently, if the implementation treats the gain as the solved matrix
$S_tK_t^\top=H_tP_{t|t-1}$, then
\begin{equation}
\label{eq:bf-score-gain-first-solve}
  \dot K_t^{(i)\top}
  =
  \mathcal{S}_{S_t}
  \left(
    \dot H_t^{(i)}P_{t|t-1}
    +H_t\dot P_{t|t-1}^{(i)}
    -\dot S_t^{(i)}K_t^\top
  \right).
\end{equation}
The filtered mean and covariance derivatives are
\begin{align}
\label{eq:bf-score-filtered-mean-first}
  \dot m_{t|t}^{(i)}
  &=
  \dot m_{t|t-1}^{(i)}
  +\dot K_t^{(i)}v_t
  +K_t\dot v_t^{(i)},\\
\label{eq:bf-score-filtered-cov-first}
  \dot P_{t|t}^{(i)}
  &=
  -\dot K_t^{(i)}H_tP_{t|t-1}
  -K_t\dot H_t^{(i)}P_{t|t-1}
  +(I-K_tH_t)\dot P_{t|t-1}^{(i)}.
\end{align}
Equation \eqref{eq:bf-score-filtered-cov-first} differentiates the simplified
covariance update $P_{t|t}=(I-K_tH_t)P_{t|t-1}$.  A Joseph-form or square-root
backend may use a different algebraic update for numerical reasons, but it must
produce derivatives of the same filtered covariance it uses in the next
prediction step.

\section{Masks and Time-Varying Observation Blocks}
\label{sec:bf-score-masks}

For missing data, let $M_t$ select the observed rows at time $t$.  The effective
observation objects are
\[
  y_t^o=M_ty_t,\qquad
  a_t^o=M_ta_t,\qquad
  H_t^o=M_tH_t,\qquad
  R_t^o=M_tR_tM_t^\top .
\]
The score recursion above applies after replacing
$(y_t,a_t,H_t,R_t)$ by $(y_t^o,a_t^o,H_t^o,R_t^o)$.  The mask is data, not a
differentiated parameter, so $\dot M_t=0$.  If no rows are observed, the
likelihood and score contribution for that time are zero and the update is the
identity:
\[
  m_{t|t}=m_{t|t-1},\qquad
  P_{t|t}=P_{t|t-1},
\]
with the same equalities for first derivatives.  This convention ensures that
the derivative path corresponds to the same masked likelihood value used by the
filter.

\section{Evidence Status}
\label{sec:bf-score-evidence-status}

The current MathDevMCP AST audit identifies Cholesky/solve, prediction, update,
quadratic-form, and derivative-recursion structure in that file, but it does
not by itself certify the algebra.  The stronger evidence is the
MacroFinance test suite comparing solve, QR square-root, autodiff, and finite
difference results on controlled LGSSMs.
